\section{Selmer ranks and local conventions}\label{sec:preliminaries}

We first record the arithmetic implications used to pass between curves
and to interpret the final analytic nonvanishing. They apply to every
elliptic curve over $\Q$.

Throughout, we use Mordell--Weil finite generation and the positive
definiteness of the canonical height modulo torsion over number fields
\citep[Chapter~IV, finite-basis theorem, Theorem~4.7,
Remark~4.12, and Theorem~6.1]{MilneEC2021}.

\begin{lemma}[Ranks and the forward implication]\label{pr:ranks}
For an elliptic curve $A/\Q$ there is an exact sequence
\[
 0\longrightarrow A(\Q)\otimes\Qtwo/\Ztwo
 \longrightarrow\Sel_{2^\infty}(A/\Q)
 \longrightarrow\Sha(A/\Q)[2^\infty]\longrightarrow0.
\]
In particular
\[
 c_2(A)=r(A)+\corank_{\Ztwo}\Sha(A/\Q)[2^\infty].
\]
If $a(A)\le1$, then $r(A)=a(A)$, the whole group $\Sha(A/\Q)$ is
finite, and $c_2(A)=a(A)$.
\end{lemma}

\begin{proof}
The finite Kummer exact sequences are
\[
 0\longrightarrow A(\Q)/2^nA(\Q)
 \longrightarrow\Sel_{2^n}(A/\Q)
 \longrightarrow\Sha(A/\Q)[2^n]\longrightarrow0.
\]
Their direct limit gives the displayed sequence; taking coranks gives
the identity. The last assertion is the Gross--Zagier--Kolyvagin
forward theorem \citep{GrossZagier1986,Kolyvagin1990}; see also its
unrestricted formulation in \citet[\S1]{Skinner2020}.
\end{proof}

\begin{lemma}[Parity and change of curve]\label{pr:invariance}
\label{prop:selmer-parity}
For every elliptic curve $A/\Q$ with root number $\varepsilon(A)$,
\[
 (-1)^{c_2(A)}=\varepsilon(A).
\]
A rational isogeny preserves $a(A)$, $r(A)$ and $c_2(A)$.
For a quadratic field $K=\Q(\sqrt{k})$,
\[
 \corank_{\Ztwo}\Sel_{2^\infty}(A/K)
       =c_2(A)+c_2(A^{(k)}).
\]
\end{lemma}

\begin{proof}
The parity statement is the $p=2$ case of
\citet[Theorem~1.4, equivalently Theorem~4.19]{DokchitserDokchitser2010}.
At $p=2$ this is Monsky's theorem \citep{Monsky1996}.
Isogenies and their duals induce maps respecting every local Kummer
condition; their composites are multiplication by their fixed degree.
After rationalization these maps preserve ranks and coranks. Their
$L$-functions agree by isogeny invariance.

Restriction to $K$ and the two conjugation projections $1+c$ and $1-c$
give maps between $\Sel_{2^\infty}(A/K)$ and the sum of the Selmer
groups of $A$ and $A^{(k)}$ over $\Q$, with composites multiplication
by $2$. They respect local Kummer images, so rationalization gives the
corank identity. An integral direct sum is not needed.
\end{proof}

\begin{lemma}[The residual alternatives]\label{pr:residual-cases}
The representation $E[2]$ is reducible over $\F_2$ if and only if
$E(\Q)[2]\ne0$. In that case it is an extension of two trivial
$\F_2$-representations. Otherwise its image is $C_3$ or $S_3$ in
$\GL_2(\F_2)$. Quadratic twisting leaves this representation unchanged.
\end{lemma}

\begin{proof}
An invariant line over $\F_2$ has a unique nonzero vector, which must
be fixed. Conversely a rational point of order two spans an invariant
line. Both that line and the one-dimensional quotient are trivial,
since $\F_2^\times=\{1\}$. The irreducible subgroups of
$\GL_2(\F_2)\simeq S_3$ are $C_3$ and $S_3$. Finally a quadratic
character reduces to $1$ modulo $2$.
\end{proof}

The distinction here is irreducibility over $\F_2$, rather than
absolute irreducibility: the cyclic cubic case is included.
In the rational-torsion case we may first replace the curve by a
rationally isogenous one with full rational two-torsion, if such a
curve exists. Otherwise, for any chosen rational two-isogeny
$E\to E'$, both $\Q(E[2])$ and $\Q(E'[2])$ are quadratic.
Lemma~\ref{pr:invariance} justifies this choice before any constructions.

\subsection{Finite lengths and twisting data}

All valuations above $2$ are normalized by $v(2)=1$, with $v(0)=+\infty$.
For a finite $\Ztwo$-module $M$, its length is $\log_2\#M$.
The corank of a cofinitely generated discrete $\Ztwo$-module is the
rank of its Pontryagin dual. Its quotient by the maximal divisible
subgroup is finite. This finite quotient must be distinguished from
finiteness of the original group.

A fixed support $S$ is a finite set of primes containing $2$ and
all primes of bad reduction; it is enlarged explicitly when necessary.
Write $N$ for the conductor of $E$.
A local filter specifies a real sign and unit squareclasses at $S$
for a variable odd fundamental discriminant prime to $S$.
The empty discriminant $h=1$ denotes the untwisted curve.
We also write $E^{(h)}$ for twisting by a rational squareclass.
For an odd prime $p$, the signed prime discriminant is
$p^*=(-1/p)p$. Products of distinct $p^*$ are odd fundamental
discriminants, with the empty product interpreted as $1$.

For an odd fundamental discriminant $h$ prime to $2N$, the
coprime-twist root-number formula is
\[
 \varepsilon(E^{(h)})=\varepsilon(E)\chi_h(-N)
\]
\citep[Section~1]{RadziwillSoundararajan2015}.
Thus matching unit squareclasses at the primes dividing $2N$ and
opposite real signs give opposite functional signs: the values of
$\chi_h(N)$ agree, whereas $\chi_h(-1)=\operatorname{sgn}(h)$.
This verifies the sign condition when choosing the auxiliary partners
below.

Put $T=T_2(E)$ and $V=E[2]$. Frobenius is arithmetic, and
$a_p(E)=p+1-\#E(\F_p)$ at a good prime. For $h$ prime to $S$, set
\begin{equation}\label{eq:h:weights}
 t_p=\dim_{\F_2}V^{\Fr_p},\qquad
 w(h)=\frac12\sum_{p\mid |h|}t_p.
\end{equation}
The residual identity, a transposition and a rotation of order three
have $t_p=2,1,0$, respectively. We call these split, simple and root
primes; their weights are $1,1/2,0$. Only types present in the
residual image are used. Quadratic residue symbols written additively
take the values $0,1\in\F_2$, corresponding to $+1,-1$.

Whenever its group is finite, write
\[
 s(h)=\length_{\Ztwo}\Sha(E^{(h)}/\Q)[2^\infty].
\]
This includes every twist already proved to have analytic rank zero
or one. It also includes a base of Selmer corank zero, since its full
$2$-primary Selmer group is then finite. We make no whole-Sha
finiteness assertion at such a base until analytic nonvanishing has
been established.

Let $\Omega_E^+$ be a positive real period and let $\Omega_E^-$ be
the positive magnitude of a nonzero anti-invariant period. Put
\begin{equation}\label{eq:h:algebraic-value}
 \mathcal L(h)=
 \frac{\sqrt{|h|}\,L(E^{(h)},1)}{\Omega_E^{\operatorname{sgn}(h)}}.
\end{equation}
Valuations refer to this algebraic period-normalized value at a chosen
place above $2$. Changes of the fixed period, modular parametrization
or isogeny affect valuations by bounded constants. All parametrizations
send the cusp $\infty$ to the origin.
Replacing Equation~\eqref{eq:h:algebraic-value} by the central value
divided by a N\'eron real period of $E^{(h)}$ changes a finite valuation
by $O_E(1)$: the differential changes by the usual twisting factor,
and there are only finitely many twisting classes at $2$.

Every uniform bound below specifies its fixed data. Constants may
depend on the curve, a chosen fixed twist and finite local preparation;
they may not depend on the subsequent cube dimension or the number of
new primes. At finite precision, the sizes of primes realizing a
specified configuration may depend on that entire configuration.
A limiting construction will only be used through its finite
realizations.

\begin{lemma}[Finite descent and corank]\label{pr:finite-descent}
Let $F$ be the quotient of $\Sha(E/\Q)[2^\infty]$ by its maximal
divisible subgroup. Then
\[
 d_2(E)=c_2(E)+\dim_{\F_2}F[2],
 \qquad \dim_{\F_2}F[2]\text{ is even}.
\]
Thus $d_2(E)\in\{0,1\}$ implies $c_2(E)=d_2(E)$ and $F=0$.
\end{lemma}

\begin{proof}
The finite Kummer sequence at $n=1$ gives
$d_2(E)=r(E)+\dim_{\F_2}\Sha(E/\Q)[2]$.
The divisible subgroup is a direct sum of copies of
$\Qtwo/\Ztwo$ and splits as an abstract abelian subgroup.
Its $2$-torsion dimension is its corank. Lemma~\ref{pr:ranks}
therefore gives the displayed identity.
The Cassels--Tate pairing on $F$ is perfect. For an elliptic curve,
the rational divisor $[O]$ defining its principal polarization makes
the pairing alternating
\citep[Chapter~I, Remark~6.10(b) and Theorem~6.13(a)]{MilneADT2006}.
A finite abelian $2$-group with such a pairing is an orthogonal sum
of pairs of cyclic groups of equal order: choose an element of
maximal order, pair it with an element detecting that order, split
off their nondegenerate plane, and continue. Hence its $2$-rank
is even. A finite $2$-group with zero $2$-rank is zero, proving
the last assertion.
\end{proof}
